% Limitations, Xingjian Liu. DRAFT, 2026-11-12: written from the MiMo-VL-7B-RL runs; the \tbd parts
% wait for the other three models and the teammates' sections (budget curves, error taxonomy).
% Needs xcolor (for \tbd) and the proposal's references.bib.
\providecommand{\tbd}{\textcolor{red}{TBD}}

\section*{Limitations}

\paragraph{Greedy decoding, one run.} Every number comes from one greedy run per model and probe. We
do not know how much the effects move under sampling, and a different decoding setup could change
how often the thinking models loop.

\paragraph{Token-limit cuts.} MiMo-VL-7B-RL reaches the 8,192-token limit in some calls, mostly by
repeating itself while thinking (on Unseen Adversarial items, 8.5\% of Default calls for B and 4.5\%
of one-shot teacher-forcing calls). The shared protocol scores a cut call as wrong, and a cut call
can never be A, so the cut rate alone can shift a difference between conditions. We therefore repeat
every comparison on the images where neither call was cut. The main effects keep their size there
(e.g.\ $+11.4$ instead of $+13.2$ points of B=A for one-shot teacher forcing on AAAAB), but small
differences in accuracy should be read with this in mind. \tbd{} (cut rates of the other models.)

\paragraph{The probes change more than the labels.} One-shot teacher forcing and the counterfactual
prefix also change the question: B is asked after a list of given labels instead of inside one
five-object answer. The dose series compares against $n=0$, which keeps that format and differs only
in the labels, but $n=0$ still differs from Default by its \texttt{obj5:} cue. The dose series also
gives the labels in box order (objects $1,\dots,n$, or $2,\dots,n+1$ in BAAAA), so the number of
labels and which boxes are labelled change together. Because MiMo-RL rewrites all five labels in its
answer, the dose changes the labels in the prompt, not how many labels come before B in the answer. Under CF the wrong class and the new question
come together, so CF alone does not show whether the drift to A comes from the model's own earlier
labels or from the question.

\paragraph{Labels in the answer come after the reasoning.} The placement probe writes the labels at
the start of the model's answer, after its reasoning has finished; we did not edit the reasoning
itself. It therefore tests whether the answer line carries inertia, not whether labels inside the
reasoning would. We score the first label written after the forced line. When the model instead
starts its own list on a new line (4--13\% of MiMo-RL's answers with true labels, 28--50\% with
counterfactual ones), that first label counts as wrong and not A, and we report the final label
separately.

\paragraph{Parsing free text.} Labels are read with a tolerant parser that maps free text to the
candidate list, and with ROPE's official strict parser. The two disagree often on MiMo-RL, which wraps answers
in \texttt{\textbackslash boxed\{\}} and re-checks its answers in prose; we report both and checked
48 core forcing rows by hand, in all of which the tolerant label matched the answer. \tbd{} (hand
check of the finer probes and of the other models.)

\paragraph{Sample size and scope.} Adversarial ROPE has 502 images (168 Seen and 334 Unseen per
order). Several effects are significant on Unseen but not on the smaller Seen split, although they
point the same way. C near / far uses the similarity of class-name embeddings
(all-mpnet-base-v2), not visual similarity, and takes A as the majority gold class. We study four
7--8B models from two families on one benchmark, so the results may not hold for larger models,
other families or other multi-object tasks. \tbd{}
